% ============================================================
% Part III §III.0 — Introduction: Attention, the Digital Habitat,
% and the Reanalysis Stance — DRAFT v1 (FCDP v2)
% Pack: plans/packs/part-iii-0-pack-v1.md · Plan: part-iii-0-D1-v1.md
% Skeleton: D1.5 run inline (no SURFACE-TO-USER items; C1–C4 dispositions ANSWER/CONCEDE-AND-LIMIT)
% Quotes enter by «Qnn» substitution only (part-iii-0-quotes.json)
% ============================================================

\section{Introduction: Attention, the Digital Habitat, and the Reanalysis Stance}

Where the account so far has moved from the inside outward---from the motion of the soul to the designed worlds that solicit it---the present Part moves back: from the frame to the laboratory, from the mechanism to the measurements, from the claim that a designed world runs or stalls the actualization of desire to three studies and one living course in which students were watched, asked, and instrumented while a designed world did exactly that. The experiments are the author's own. The virtual learning environment under study was built for an undergraduate biomedical engineering course; the studies that instrumented it---one physiological, one oculometric, one phenomenological---were conducted and published with its builder's hand in each; and the course itself, run fully online and at scale, has now yielded three terms of student testimony. What follows is therefore not a new experiment but a reanalysis: the published findings and the primary record are read again, through the frame this dissertation has built, so that the data yield the conclusions and the frame supplies what data alone cannot---the how and the why of their arising.

Three goals govern the Part. The first is boredom: to analyze the student's boredom phenomenologically and rhetorical-ontologically, so that what the instruments recorded as signal becomes legible as attunement. The second is engagement: to assess what the surveys and studies called engagement with the dissertation's own methodology, and to say what that word was measuring when it measured anything at all. The third is design: to show that a frame which explains why virtual worlds bore can direct the making of ones that do not---that diagnosis, carried through, becomes a design discipline. Boredom explained, engagement assessed, design instructed: the three goals are one movement performed three times, each at a different depth of the same chain.

One fact about the corpus must stand at the door, because everything after passes through it. Between the pilot studies and the deployment, the course moved: virtual reality remained available, but a full-sized class, run wholly online and asynchronously, could not be scaled into headsets, and so the environment met its students as a three-dimensional world on a two-dimensional screen. The move was not a fall from the true experiment. It was the discovery of the actual one.

Why should engagement with digital artifacts be analyzed at all? The question deserves its bluntness, and the answer begins a century before the screen. William James, whom Gloria Mark sets at the head of her account of the modern attention span, held that attending is not a lamp the mind plays over a finished world but the very act by which a life takes shape: ``«Q1»''«Q1+». Experience is not what happens; experience is what is attended to while it happens. Mark carries the claim into the present tense of every pocket: walking through a garden with a phone in hand, she observes, it is the exchange of messages that survives as the walk's record---``«Q2»''«Q2+». The garden was present; it was not attended; it never properly entered. What holds the attention holds, in the end, the life.

And attention is held, now, by screens---held briefly, and held often. Mark's own logging studies, replicated across teams and workplaces, find that ``«Q3»''«Q3+»; by 2016 the median stretch had fallen to forty seconds. Whatever else the digital artifact is---tool, medium, world---it is the place where attending now factually lives, in tenancies measured in seconds and ended by a switch. To study human engagement with digital artifacts is not, then, a regional curiosity of educational technology. It is the study of attention in its native habitat, conducted where the constitution of experience is decided several hundred times in a working day.

An objection stirs: the Jamesian warrant proves too much, for by it every medium---book, lecture, window---would equally constitute experience, and nothing particular would follow for the screen. The concentration answers it. The claim is not that screens alone shape what enters a life but that they are where the shaping now happens with measured regularity and measured brevity; the forty-seven seconds are not a metaphysical thesis but a census. Where attention lives is an empirical question, and it has received an empirical answer.

The dissertation's frame says the same thing in its own tongue. What is attended is what enters the chain---becomes \textit{phantasma}, takes \textit{doxa}, issues in act---and what enters the chain, repeated, sediments into the \textit{hexis} that a student carries out of the course. Attention is not the antechamber of learning; it is learning's first act. A world that cannot hold attention cannot form anyone.

How attention is studied has meanwhile turned on its own axis. Where an earlier psychology asked and tested, the present science senses: ``«Q4»''«Q4+». Mark reports the turn as an accomplished fact of the field, and the corpus of this Part enacted it in miniature before the frame ever arrived. Across the three studies the heart-rate measures were dropped as uninformative; the electroencephalogram, too noisy in practice, was demoted to corroboration; the eye, tracked at a hundred and twenty hertz, was elevated to the primary instrument; and the third study turned to the phenomenological interview. The program walked, on its own methodological legs, from asking toward sensing and from sensing toward the description of lived experience. The walk matters because it was not driven by theory. It was driven by what the instruments could and could not see---and what they could not see singly is precisely what this Part will name.

The same instrumentation has a mirror. Industrially deployed, the sensors that might diagnose a stalled attention are used instead to seize it: ``«Q5»''«Q5+». The notification engineered for the bored is the dark twin of the lesson designed for the engaged---both know the stalled chain; one exploits the stall, the other would repair it. That the knowledge is symmetrical is the strongest reason to hold it: a pedagogy that declines to understand why elements bore does not thereby keep its hands clean; it merely leaves the understanding to those who monetize it. The mirror will return at the Part's end, where the norms of design are at stake. Here it sets the stakes.

A second objection must be met, for it cuts at the corpus itself. If surveys and tests are superseded by sensors, what standing have the three terms of survey testimony on which the deployment analysis will lean? The concession is real, and it is limited: no single channel---sensed or asked---ever suffices for the phenomenon at issue, and the frame's central methodological rule will be precisely that the deepest form of boredom shows itself only where two channels disagree. The survey is one voice, indispensable and insufficient; the instruments are another; the analysis listens for their divergence.

One misreading must be blocked before the corpus is opened, because the vocabulary invites it. To say that a designed environment runs or stalls the actualization of desire sounds, to a certain ear, like the oldest behaviorism---``«Q6»''«Q6+». The resemblance is surface. In the frame advanced here the designer and the participant divide the causes of the act as the producing art divides them from the using art: the designer furnishes the matter and the form---the rooms, the images, the affordances---and by furnishing them occasions the act; the efficient and the final cause remain the participant's own, or there is no act at all, only motion. A world conditions nothing; it offers. What the participant does with the offer---whether the \textit{phantasma} is taken up, whether \textit{doxa} commits, whether anything is done---is the participant's actualization, not the environment's output. Cognition, far from being a fallacy, is the very site where the offer is accepted or refused; the whole analysis to follow is an anatomy of that acceptance and that refusal. So when a later section says that a world gave its students nothing to do, the grammar of the complaint should be heard exactly: not that the world failed to produce behavior, but that it withheld the conditions under which a student could complete an act of one's own. The designer's power is real and bounded: real, because the offer shapes everything; bounded, because the offer is all it is. Design, on these terms, is rhetoric---and rhetoric persuades agents; it does not program mechanisms.

The corpus, then. At one pole stand the laboratory and the pilot. Between 2022 and 2024 three published studies instrumented the environment's early life: a physiological pilot that watched three participants' brains and hearts across three seventeen-minute films; an eye-tracking study that followed twelve participants' gazes at a hundred and twenty hertz; a phenomenological evaluation that asked and interviewed a pilot cohort while four formats of the same clinical footage competed for their preference. Around them stands the raw record itself, newly assembled for this Part: eight participants' electroencephalograms, headset telemetry, session recordings, and---most valuable of all, it will emerge---their written answers, film by film, to seven blunt questions, among them how long each film felt. At the other pole stands the deployment: three terms of the live course, two hundred forty-one students, eight hundred seventy-three written answers, a world reached through a downloadable application on the student's own computer or, for nearly half of every term with remarkable constancy, not reached at all in favor of the plain videos. Between the poles falls the move itself: presence spoke for the headset, feasibility spoke for the screen, and the course, like most of its students, followed feasibility.

The record, read cold, already leans toward the frame. Its instruments disagreed---and the disagreement is the door. On the same clinical film the earlier study's electroencephalogram read a profile that resembled boredom, while the later study's gaze-variance read engagement with statistical confidence; the participants meanwhile called themselves interested, and felt seventeen minutes as twenty-five or thirty. No bipolar scale can hold that record. Each instrument, alone, tells a different story; and the program, to its credit, kept the contradiction rather than smoothing it. The frame will read the contradiction in one move---a surface occupied, a depth left empty---and the reading will organize everything that follows: the laboratory's forms of boredom, the deployment's flights and frictions, the design levers that answer them.

The route runs so: first the method of reanalysis, stated with its limits; then the laboratory, where the forms of boredom acquire empirical faces; then the deployment, where a designed world met two hundred forty-one students and mostly showed them things; then the synthesis, where the corpus is read as one and the question why elements bore receives its answer; last the return to design, where the lesson that to design a world is to design the conditions under which a chain runs or stalls is brought home to the classroom with the corpus's own evidence in hand. The data are already in hand. What they awaited was the frame.

% ============================================================
% FOOTER (FCDP provenance)
% Material: ALL NEW (greenfield; no foundation text; P7 = N/A).
% Retained: n/a. Corrections: n/a.
% Omitted: E8 (Mark college-students/children passages) — pages unverified, per grade rule.
% Quotes used: Q1–Q6 (all ASSIGNED, all used). Q5 citation carries unverified-page flag.
% Counterargument dispositions honored: C1 ANSWER (M2 ¶3); C2 CONCEDE-AND-LIMIT (M3 ¶3);
%   C3 ANSWER (M4); C4 ANSWER (M5 ¶2 — the kept contradiction as proof of need).
% ============================================================
